\lysh{37171}{2}

\begin{alist}
\item Είναι:
\begin{align*}
\alpha^2\beta + \alpha\beta^2 = -30 \ &\Leftrightarrow
\alpha\beta(\alpha + \beta) = -30 \ \Leftrightarrow \\
\alpha\beta \cdot 2 = -30 \ &\Leftrightarrow
\alpha\beta = -15
\end{align*}

\item Η ζητούμενη εξίσωση μπορεί να είναι της μορφής $x^2 - Sx + P = 0$, με $S = \alpha + \beta = 2$ και $P = \alpha\beta = -15$.
Τελικά, μία ζητούμενη εξίσωση είναι η $x^2 - 2x - 15 = 0$.

Το τριώνυμο $x^2 - 2x - 15$ έχει διακρίνουσα $\Delta = \beta^2 - 4\alpha\gamma = (-2)^2 - 4 \cdot 1 \cdot (-15) = 4 + 60 = 64 > 0$.
Οι ρίζες της εξίσωσης είναι:
\[x_{1,2} = \dfrac{-\beta \pm \sqrt{\Delta}}{2\alpha} = \dfrac{-(-2) \pm \sqrt{64}}{2 \cdot 1} = \dfrac{2 \pm 8}{2} = \begin{cases} 5 \\ -3 \end{cases}\]

Άρα είναι $\alpha = 5$ και $\beta = -3$ ή $\alpha = -3$ και $\beta = 5$.
\end{alist}

